\documentclass[11pt, draft]{article}
\topmargin=0mm
\oddsidemargin=0mm
\textwidth170mm
\textheight230mm
\begin{document}
%\pagestyle{empty}
\language=0
%\hyphenation{}

\begin{center}
{\bf
COMPUTATION OF JANET BASES. I.MONOMIAL BASES}\\
\vspace {0.3cm}
{Vladimir Gerdt$^1$, Yuri Blinkov$^2$, Denis Yanovich$^1$}\\
{\it \small$^1$Laboratory of Computing techniques and Automation, Joint Institute for Nuclear
Research, 141980 Dubna, Russia e-mail: gerdt@jinr.ru\ \ yan@cv.jinr.ru \\
$^2$Department of Mathematics and Mechanics, Saratov University, 410071 Saratov, Russia \\
e-mail: BlinkovUA@info.sgu.ru}
\end{center}
\vspace {1.0cm}
In this talk we present algorithms and their implementation in C
for construction of monomial Janet bases. As data structures for
finite monomial sets certain trees called Janet trees are selected.
We describe an algorithm for construction of a minimal Janet basis for the
ideal generated by a finite number of monomials which is the
specific form of that in paper~\cite{GB2,G98} for Janet division.
As subalgorithms the algorithm contains those for fast search of Janet
divisors in the given Janet tree, and for insertion of irreducible
nonmultiplicative prolongation into this tree. The presented algorithms
implemented in C. We give some examples illustrating practical
efficiency of the algorithms and implementation. The algorithms
play also an important role in computation of polynomial Janet bases.
\vskip 0.5cm

\begin{thebibliography}{10}

\bibitem{GB2} Gerdt V.P., Blinkov A.Yu.{\em Minimal Involutive Bases}.
 Mathematics and Computers in Simulation 45 (1998) pp.543-560. \\
 URL: http://arXiv.org/abs/math.RA/9912029.
\bibitem{G98} Gerdt V.P. {\em Involutive Division Technique: Some
 Generalizations and Optimizations}. Zapiski Nauchnykh Seminarov POMI (St.Petersburg)
 {\bf 258} (1999) 185-206. To appear in {\em J. Math. Sci.} {\bf 258}
(2000).\\  URL: http://arXiv.org/abs/math.RA/9912030.

\end{thebibliography}

\end{document}
